% =====================================================================
%  Quickstart de LaTeX — tu primer documento
%  Suite de documentos LaTeX · Nivel 1 (Empezar) · clase article.
%  Compilar: latexmk -pdf quickstart.tex
% =====================================================================
\documentclass[11pt,a4paper]{article}

% --- Base estándar (ver ESTANDARES_LATEX.md) ---
\usepackage[utf8]{inputenc}
\usepackage[T1]{fontenc}
\usepackage{lmodern}
\usepackage{microtype}
\usepackage[spanish,es-noshorthands]{babel}
\usepackage{csquotes}
\usepackage{amsmath}
\usepackage[margin=2.4cm]{geometry}
\usepackage{enumitem}
\usepackage{xcolor}
\usepackage{listings}
\usepackage{hyperref}
\hypersetup{unicode, pdftitle={Quickstart de LaTeX}, pdfauthor={Rodrigo Martín Vidal Galán},
            colorlinks=true, linkcolor=blue!50!black, urlcolor=blue!50!black}
\usepackage[spanish,capitalise,nameinlink]{cleveref}  % después de hyperref

% --- Estilo de código ---
\definecolor{codeback}{RGB}{245,247,250}
\definecolor{codekw}{RGB}{20,60,120}
\definecolor{codecom}{RGB}{120,120,120}
\lstset{
  language=[LaTeX]TeX,
  basicstyle=\ttfamily\small,
  keywordstyle=\color{codekw}\bfseries,
  commentstyle=\color{codecom}\itshape,
  backgroundcolor=\color{codeback},
  frame=single, rulecolor=\color{codekw!30},
  breaklines=true, showstringspaces=false, columns=fullflexible,
  xleftmargin=10pt, xrightmargin=4pt, framesep=4pt,
  extendedchars=true,
  literate={á}{{\'a}}1 {é}{{\'e}}1 {í}{{\'i}}1 {ó}{{\'o}}1 {ú}{{\'u}}1
           {ñ}{{\~n}}1 {¿}{{?`}}1 {¡}{{!`}}1 {°}{{\textdegree}}1,
}
\newcommand{\C}[1]{\texttt{\textbackslash #1}}   % imprime el nombre de un comando

% --- Cajas de ayuda ---
\newcommand{\pagina}[1]{%
  \begin{center}\fbox{\begin{minipage}{0.82\linewidth}\small #1\end{minipage}}\end{center}}
\newenvironment{tip}{\par\smallskip\noindent\begin{center}%
  \begin{minipage}{0.94\linewidth}\small\color{codekw!85!black}\rule{\linewidth}{0.4pt}\par\nobreak
  \textbf{Tip.}\enspace}{\par\rule{\linewidth}{0.4pt}\end{minipage}\end{center}\smallskip}

\title{\textbf{Quickstart de \LaTeX}\\[2pt]\large Tu primer documento, de cero}
\author{Rodrigo Martín Vidal Galán}
\date{Junio de 2026}

\begin{document}
\maketitle

\noindent\textbf{¿Para quién es esto?} Para quien \emph{nunca} usó \LaTeX{} (o lo
usó muy por arriba). Es tu \textbf{primer encuentro} con la herramienta: arrancamos
de cero, instalamos, escribimos, compilamos y armamos un documento de verdad,
explicando el \emph{porqué} de cada cosa y anticipando las dudas típicas del
principio. No profundiza en ningún tema (para eso están los demás documentos de la
suite); el objetivo es que termines \textbf{cómodo y compilando}.

\newpage
\tableofcontents
\newpage

% =====================================================================
\section{¿Qué es LaTeX y en qué se diferencia de Word?}
% =====================================================================
\LaTeX{} es un \textbf{sistema de composición de documentos}: vos escribís texto
plano con \emph{instrucciones} (comandos que empiezan con \texttt{\textbackslash})
y \LaTeX{} se encarga de la \textbf{apariencia} (tipografía, espaciado, fórmulas,
numeración).

La diferencia clave con Word es que \LaTeX{} \textbf{no es \enquote{lo que ves es
lo que obtenés}} (no WYSIWYG): mientras escribís, ves \emph{código}, no el
resultado final; para ver cómo queda, \textbf{compilás} (un paso que convierte tu
texto en un PDF). A cambio, \LaTeX{} \textbf{separa el contenido de la
presentación}: vos decís \emph{qué} es cada cosa (\enquote{esto es un título},
\enquote{esto es una fórmula}) y \LaTeX{} decide cómo se ve, de forma consistente
en todo el documento. Por eso es el estándar para matemática y textos técnicos.

\begin{tip}
Al principio se trabaja \enquote{a ciegas}: escribís un poco y compilás para ver.
Es normal que cueste los primeros días; no hace falta memorizar nada (más adelante
vemos cómo buscar lo que necesites).
\end{tip}

% =====================================================================
\section{Instalación}
% =====================================================================
Tenés dos caminos. Si querés evitar instalar, empezá por el A.

\paragraph{Opción A — Sin instalar nada (Overleaf).}
\begin{enumerate}[noitemsep]
  \item Entrá a \texttt{overleaf.com} y creá una cuenta gratis.
  \item \emph{New Project} $\to$ \emph{Blank Project}.
  \item Ya tenés un editor que \textbf{compila en el navegador}. Ideal para arrancar.
\end{enumerate}

\paragraph{Opción B — Instalar en tu computadora.} Necesitás \emph{dos} cosas:
\begin{enumerate}[noitemsep]
  \item Una \textbf{distribución} de \LaTeX{} (el motor + los paquetes): bajá
        \texttt{MiKTeX} de \texttt{miktex.org} (Windows), o \texttt{TeX~Live}
        (multiplataforma) / \texttt{MacTeX} (macOS). Corré el instalador.
  \item Un \textbf{editor}: \texttt{TeXstudio} (\texttt{texstudio.org}) o
        \texttt{VS~Code} con la extensión \emph{LaTeX Workshop}.
\end{enumerate}
Después abrís el editor, creás un archivo \texttt{.tex} y compilás (sección
siguiente).

\begin{tip}
La \textbf{primera} vez que compiles, MiKTeX te va a ofrecer instalar los paquetes
que falten: aceptá. (Para el detalle de distribuciones y sus problemáticas, ver el
documento \emph{Flujo e instalación}.)
\end{tip}

% =====================================================================
\section{Tu primer documento (\enquote{hola mundo})}
% =====================================================================
Creá un archivo \texttt{hola.tex} con \emph{exactamente} esto:
\begin{lstlisting}
\documentclass{article}     % tipo de documento
\begin{document}            % --- empieza el cuerpo ---
Hola, mundo.
\end{document}              % --- termina el cuerpo ---
\end{lstlisting}
Al compilarlo (próxima sección) obtenés una página con:
\pagina{Hola, mundo.}

\noindent\textbf{Las tres partes de todo \texttt{.tex}:}
\begin{enumerate}[noitemsep]
  \item \C{documentclass\{...\}}: el \textbf{tipo} de documento
        (\texttt{article} para artículos/TPs, \texttt{report} con capítulos,
        \texttt{beamer} para presentaciones\dots).
  \item \textbf{Preámbulo}: todo lo que va \emph{entre} el \C{documentclass} y el
        \C{begin\{document\}}. Acá se cargan paquetes y se configuran cosas (más
        abajo lo completamos).
  \item \textbf{Cuerpo}: lo que va entre \C{begin\{document\}} y
        \C{end\{document\}}. Es lo que realmente aparece en el PDF.
\end{enumerate}
Las líneas que empiezan con \texttt{\%} son \textbf{comentarios}: \LaTeX{} las
ignora (sirven para anotarte cosas).

% =====================================================================
\section{Compilar (y qué pasa cuando compilás)}
% =====================================================================
\emph{Compilar} es convertir tu \texttt{.tex} en un \texttt{.pdf}:
\[ \texttt{archivo.tex} \;\xrightarrow{\;\text{compilar}\;}\; \texttt{archivo.pdf} \]
\begin{itemize}[noitemsep]
  \item \textbf{En el editor:} apretá \emph{Compilar} / \emph{Build} (en TeXstudio,
        la tecla \texttt{F5}; en Overleaf, \emph{Recompile}).
  \item \textbf{En la terminal:} \texttt{latexmk -pdf archivo.tex} (resuelve solo
        todas las pasadas necesarias).
\end{itemize}
Al compilar, además del \texttt{.pdf} aparecen \textbf{archivos auxiliares}
(\texttt{.aux}, \texttt{.log}, \texttt{.toc}\dots): son normales y se pueden
ignorar; el \texttt{.log} guarda el detalle de errores. Si hay un \textbf{error}, la
compilación se detiene y el editor te marca la línea (ver \cref{sec:errores}).

\begin{tip}
Cuando uses \textbf{referencias} o \textbf{índice}, hay que compilar \textbf{dos
veces} (la primera \enquote{anota} los números, la segunda los coloca).
\texttt{latexmk} lo hace solo.
\end{tip}

El ciclo de trabajo es: \textbf{escribir $\to$ compilar $\to$ mirar el PDF $\to$
corregir $\to$ repetir}.

% =====================================================================
\section{Los ladrillos: comandos, entornos y paquetes}
% =====================================================================
Casi todo en \LaTeX{} es una de estas tres cosas. Entenderlas te ahorra el 80\,\%
de la confusión inicial.

\paragraph{Comando.} Empieza con \texttt{\textbackslash} y puede llevar
argumentos \textbf{obligatorios} entre llaves \texttt{\{\}} y \textbf{opcionales}
entre corchetes \texttt{[]}:
\begin{lstlisting}
\textbf{en negrita}        % un argumento obligatorio
\textcolor{red}{rojo}      % dos argumentos
\rule[2pt]{2cm}{1pt}       % uno opcional [2pt] y dos obligatorios
\end{lstlisting}

\paragraph{Entorno.} Un \textbf{bloque} con apertura y cierre, que afecta a todo
lo que está adentro:
\begin{lstlisting}
\begin{center}
  esto queda centrado
\end{center}
\end{lstlisting}
Todo \C{begin\{X\}} debe tener su \C{end\{X\}}. Olvidarlo es un error clásico.

\paragraph{Paquete.} Una \textbf{extensión} que cargás en el preámbulo con
\C{usepackage} para tener más comandos disponibles:
\begin{lstlisting}
\usepackage{graphicx}   % habilita \includegraphics para poner imágenes
\usepackage{amsmath}    % matematica avanzada
\end{lstlisting}

\begin{tip}
Muchos errores de \enquote{comando indefinido} se arreglan cargando el paquete
correspondiente en el preámbulo.
\end{tip}

% =====================================================================
\section{Reglas básicas del texto (lo que sorprende al principio)}
% =====================================================================
\begin{itemize}[itemsep=3pt]
  \item \textbf{Espacios:} varios espacios seguidos cuentan como \textbf{uno}
        solo. Para forzar más espacio usá comandos (\C{quad}, \C{hspace\{1cm\}}).
  \item \textbf{Párrafos:} una \textbf{línea en blanco} empieza un párrafo nuevo.
        Un solo \emph{Enter} (sin línea en blanco) \textbf{no} corta el párrafo.
  \item \textbf{Caracteres especiales:} estos diez tienen significado propio y no
        se escriben directo: \verb|% $ & # _ { } ~ ^ \|. Para que aparezcan
        literales: \verb|\% \$ \& \# \_ \{ \}| (y \C{textbackslash} para la
        \enquote{\textbackslash}).
  \item \textbf{Comillas:} no uses las rectas del teclado. Escribí
        \verb|`así'| (simples) o \verb|``así''| (dobles) —o mejor
        \C{enquote\{así\}} (del paquete \texttt{csquotes})—.
  \item \textbf{Guiones:} \verb|-| (guion), \verb|--| (rango: \enquote{8--12}),
        \verb|---| (raya de inciso).
\end{itemize}
\noindent\textbf{Ejemplo:}
\begin{lstlisting}
El 50\% de los casos (a--b) usa \LaTeX, dijo ``el profe''.
\end{lstlisting}
produce: El 50\% de los casos (a--b) usa \LaTeX, dijo \enquote{el profe}.

% =====================================================================
\section{Texto, listas y matemática básica}
% =====================================================================
Lo más usado para empezar a escribir:
\begin{lstlisting}
Texto \textbf{en negrita} y \emph{con énfasis}.

\begin{itemize}
  \item Una viñeta
  \item Otra viñeta
\end{itemize}

\begin{enumerate}
  \item Primer paso
  \item Segundo paso
\end{enumerate}

En línea: $a^2 + b^2 = c^2$. Y centrada:
\[ x = \frac{-b \pm \sqrt{b^2 - 4ac}}{2a} \]
\end{lstlisting}
produce:
\begin{quote}
Texto \textbf{en negrita} y \emph{con énfasis}.
\begin{itemize}[noitemsep]
  \item Una viñeta
  \item Otra viñeta
\end{itemize}
\begin{enumerate}[noitemsep]
  \item Primer paso
  \item Segundo paso
\end{enumerate}
En línea: $a^2 + b^2 = c^2$. Y centrada:
\[ x = \frac{-b \pm \sqrt{b^2 - 4ac}}{2a} \]
\end{quote}
\noindent La matemática va entre \texttt{\$...\$} (en línea) o \C{[}\,\dots\,\C{]}
(centrada). Para todo lo de matemática a fondo, ver el documento \emph{Matemática}.

% =====================================================================
\section{Construí tu documento paso a paso}
% =====================================================================
Un documento de verdad se arma sumando piezas. La idea es ir \textbf{de a poco},
compilando en cada paso:
\begin{enumerate}[noitemsep]
  \item Arrancás con el \textbf{esqueleto} mínimo (\enquote{hola mundo}).
  \item Completás el \textbf{preámbulo} con idioma y fuentes
        (\texttt{babel}, \texttt{fontenc}).
  \item Agregás \textbf{título} (\C{maketitle}).
  \item Dividís en \textbf{secciones} (\C{section}, se numeran solas).
  \item Ponés el \textbf{contenido} (texto, listas, fórmulas, imágenes).
\end{enumerate}
El resultado de juntar todo:
\begin{lstlisting}
\documentclass[11pt]{article}
\usepackage[utf8]{inputenc}     % preámbulo:
\usepackage[T1]{fontenc}        %   idioma y
\usepackage[spanish]{babel}     %   fuentes

\title{Mi primer informe}
\author{Yo}
\date{\today}

\begin{document}
\maketitle

\section{Introducción}
Esto es la introducción.

\section{Desarrollo}
Y esto el desarrollo.
\end{document}
\end{lstlisting}
produce (una página):
\pagina{%
  \centering {\large\textbf{Mi primer informe}}\\[2pt]
  Yo \quad\textbullet\quad \today\\[8pt]
  \raggedright
  \textbf{1\quad Introducción}\\ Esto es la introducción.\\[4pt]
  \textbf{2\quad Desarrollo}\\ Y esto el desarrollo.}

\noindent\textbf{¿Una imagen?} Cargá \C{usepackage\{graphicx\}} y usá
\C{includegraphics[width=...]\{foto.png\}} dentro de un entorno \texttt{figure}
(lo vemos en \emph{Gráficos y figuras}).

% =====================================================================
\section{Errores comunes del día 1}\label{sec:errores}
% =====================================================================
Cuando algo falla, la compilación se \textbf{detiene} y el editor te marca la
línea. Los más típicos:
\begin{itemize}[itemsep=3pt]
  \item \texttt{Missing \$ inserted}: pusiste un símbolo matemático (como
        \verb|_| o \verb|^|) fuera del modo math. Encerralo en \texttt{\$...\$}.
  \item \texttt{Undefined control sequence}: comando mal escrito o falta un
        \C{usepackage}. Revisá la ortografía y, si es de un paquete, cargalo.
  \item \texttt{Runaway argument} / \texttt{Missing \}}: una llave o un entorno
        sin cerrar. Revisá los \C{begin}/\C{end} y las llaves.
  \item \textbf{Referencias o índice como \enquote{??}}: faltó compilar
        \textbf{dos veces} (o usá \texttt{latexmk}).
  \item \C{usepackage} \textbf{después} de \C{begin\{document\}}: los paquetes van
        en el \textbf{preámbulo}, nunca en el cuerpo.
\end{itemize}
\begin{tip}
\textbf{Leé la \emph{primera} línea de error y el número de línea}: casi siempre el
problema está ahí (o una línea antes). Los errores siguientes suelen ser
\enquote{efecto dominó} del primero.
\end{tip}

% =====================================================================
\section{Cómo seguir y pedir ayuda}
% =====================================================================
\textbf{Nadie memoriza \LaTeX:} se busca. Las tres herramientas clave:
\begin{itemize}[noitemsep]
  \item \texttt{texdoc <paquete>} en la terminal abre el \textbf{manual oficial}
        (p.\,ej.\ \texttt{texdoc amsmath}).
  \item Buscar en internet \enquote{\textbf{LaTeX cómo hacer X}} o directamente en
        \texttt{tex.stackexchange.com} (casi todo ya está preguntado).
  \item Para pedir ayuda, armá un \textbf{MWE} (\emph{minimal working example}): el
        documento más \textbf{chico} que reproduce el problema. Quien te ayuda lo
        puede compilar y ver al toque.
\end{itemize}

\paragraph{Mini-glosario.}
\begin{description}[noitemsep,leftmargin=2.4cm,style=nextline]
  \item[Comando] orden que empieza con \texttt{\textbackslash} (ej.\ \C{section}).
  \item[Entorno] bloque \C{begin\{X\}}\,\dots\,\C{end\{X\}}.
  \item[Paquete] extensión que se carga con \C{usepackage}.
  \item[Preámbulo] lo que va antes de \C{begin\{document\}}.
  \item[Compilar] convertir el \texttt{.tex} en \texttt{.pdf}.
  \item[MWE] ejemplo mínimo que reproduce un problema.
\end{description}

% =====================================================================
\section*{Próximos pasos / Ver también}
% =====================================================================
\addcontentsline{toc}{section}{Próximos pasos / Ver también}
¡Listo, ya compilás y entendés los fundamentos! Para seguir, según lo que necesites:
\begin{itemize}[noitemsep]
  \item \emph{Masterclass de \LaTeX} — una recorrida general en formato charla.
  \item \emph{Matemática} — escribir matemática a fondo.
  \item \emph{Gráficos y figuras} — imágenes y graficar funciones.
  \item \emph{Hoja de referencia} — los comandos más usados en una carilla.
  \item \emph{Estándares y buenas prácticas} — el preámbulo recomendado y cómo
        hacer las cosas bien.
  \item \texttt{Overleaf} y la guía \emph{The Not So Short Introduction to \LaTeX}.
\end{itemize}

\end{document}
